\section{Efisiensi Pembaruan Kebijakan}

H3 dan H4 diuji bersama karena keduanya berkaitan dengan siklus pembaruan kebijakan. 
H3 menguji apakah pemisahan pembaruan struktur dan parameter dapat mengurangi waktu komputasi tanpa menurunkan kualitas secara signifikan, dengan efek yang bergantung pada tingkat \textit{drift}. 
H4 menguji apakah strategi \textit{warm-start} dapat mencapai kualitas yang sekurang-kurangnya setara dengan \textit{re-mining} penuh dalam waktu lebih singkat.

Pengujian dilakukan dengan membagi \textit{log} menjadi lima jendela berurutan. 
Struktur Chow–Liu dipelajari ulang hanya ketika \textit{Jensen–Shannon divergence} melampaui ambang $\tau=0{,}01$. 
Empat kondisi dibandingkan, yaitu \textit{warm} (pembaruan struktur berbasis \textit{drift}), \textit{always} (struktur diperbarui pada setiap jendela), \textit{cold} (\textit{re-mining} penuh), dan iBOA \textit{cold}. 
Ringkasan hasil disajikan pada Tabel \ref{tab:hasil-h3h4}.

\begin{longtable}{|>{\raggedright\arraybackslash}p{2cm}|>{\raggedright\arraybackslash}p{1.5cm}|>{\raggedright\arraybackslash}p{1.5cm}|>{\raggedright\arraybackslash}p{1.5cm}|>{\raggedright\arraybackslash}p{1.5cm}|>{\raggedright\arraybackslash}p{3.5cm}|}
\caption{Hasil Pengujian H3 dan H4 pada Lima Dataset (rerata; warm/always/cold)}
\label{tab:hasil-h3h4} \\
\hline
\textbf{Dataset / skenario} & \textbf{Pemicu relearn} & \textbf{warm (s)} & \textbf{always (s)} & \textbf{cold (s)} & \textbf{Vonis} \\
\hline
\endfirsthead

\multicolumn{6}{c}{{\bfseries \tablename\ \thetable{} -- lanjutan dari halaman sebelumnya}} \\
\hline
\textbf{Dataset / skenario} & \textbf{Pemicu relearn} & \textbf{warm (s)} & \textbf{always (s)} & \textbf{cold (s)} & \textbf{Vonis} \\
\hline
\endhead

\hline
\multicolumn{6}{r}{Lanjut ke halaman berikutnya} \\
\endfoot

\hline
\endlastfoot

% ==================== DATA BARIS ====================
workforce / stasioner & 0 dari 4 & 3{,}46 & 4{,}58 & 9{,}05 & H3 TERDUKUNG ($p=6\times10^{-5}$); H4 $2{,}6\times$, F1 $+0{,}095$ \\
\hline
workforce / drift sedang & 4 dari 4 & 4{,}52 & 4{,}53 & 8{,}95 & H3 tak berlaku (struktur memang perlu); H4 $2{,}0\times$, F1 $+0{,}061$ \\
\hline
e-document / stasioner & 4 dari 4 & 22{,}3 & 22{,}3 & 28{,}2 & H3 tak dapat ditunjukkan (lantai derau $>$ $\tau$); H4 $1{,}27\times$, F1 $+0{,}102$ \\
\hline
project-management / stasioner & 0 dari 4 & 0{,}17 & 0{,}21 & 4{,}45 & H4 $25{,}9\times$, F1 $-0{,}020$ (kecil, jenuh) \\
\hline
university / stasioner & 0 dari 4 & 0{,}08 & 0{,}11 & 3{,}76 & H4 $49\times$, F1 $+0{,}007$ \\
\hline

\end{longtable}

H3 didukung secara bersyarat. 
Pada \textit{workforce} stasioner, pemicu \textit{drift} tidak aktif selama empat siklus, sedangkan kondisi \textit{always} tetap mempelajari ulang struktur pada setiap siklus. 
Akibatnya, \textit{warm-start} lebih cepat daripada \textit{always} (3,46 dibanding 4,58 detik; $p=6{,}1\times10^{-5}$) tanpa penurunan kualitas. 
Pada \textit{drift} rendah dan sedang, pemicu aktif pada setiap jendela sehingga \textit{warm} berperilaku sama dengan \textit{always}, sesuai rancangan mekanisme. 
Pada \textit{e-document}, lantai derau JSD melebihi $\tau$ sehingga pemicu terus aktif dan penghematan H3 tidak dapat ditunjukkan. 
Dengan demikian, manfaat pemisahan pembaruan struktur bergantung pada intensitas \textit{drift} dan lantai derau distribusi dataset.

H4 memperoleh dukungan kuat. 
\textit{Warm-start} mempercepat re-optimasi dibandingkan \textit{re-mining} penuh pada seluruh dataset, yaitu sekitar $2,0$–$2,6$ kali pada \textit{workforce}, $1,27$ kali pada \textit{e-document}, dan $21$–$52$ kali pada tiga dataset kecil. 
Pada dua dataset utama, \textit{warm-start} juga menghasilkan F1 lebih tinggi sebesar $0,06$–$0,17$ karena proses pembaruan mempertahankan informasi kebijakan sebelumnya dan berfokus pada pola yang belum tercakup. 
Pada tiga dataset kecil, perbedaan kualitas relatif kecil dan tidak konsisten ($-0{,}057$ hingga $+0{,}031$), sehingga bukti utama kualitas H4 berasal dari \textit{workforce} dan \textit{e-document}.

Hasil H3 dan H4 menunjukkan manfaat CoDA-EDA pada pembaruan kebijakan berkala di \textit{edge}. 
Pemicu berbasis \textit{drift} dapat menghindari pembelajaran ulang struktur ketika distribusi stabil, sedangkan \textit{warm-start} mengurangi waktu re-optimasi dibandingkan penambangan ulang dari awal. 
Dibandingkan iBOA, yang memiliki kualitas kompetitif pada penambangan statis, keunggulan CoDA-EDA lebih terlihat pada skenario pembaruan berkala. 
Bersama hasil H2, temuan ini menunjukkan bahwa kontribusi utama CoDA-EDA terletak pada efisiensi memori dan pembaruan kebijakan, bukan pada dominasi kualitas penambangan statis.